% Options for packages loaded elsewhere
\PassOptionsToPackage{unicode}{hyperref}
\PassOptionsToPackage{hyphens}{url}
\PassOptionsToPackage{dvipsnames,svgnames,x11names}{xcolor}
\documentclass[
]{ctexart}
\usepackage{xcolor}
\usepackage[margin=2.2cm]{geometry}
\usepackage{amsmath,amssymb}
\setcounter{secnumdepth}{-\maxdimen} % remove section numbering
\usepackage{iftex}
\ifPDFTeX
  \usepackage[T1]{fontenc}
  \usepackage[utf8]{inputenc}
  \usepackage{textcomp} % provide euro and other symbols
\else % if luatex or xetex
  \usepackage{unicode-math} % this also loads fontspec
  \defaultfontfeatures{Scale=MatchLowercase}
  \defaultfontfeatures[\rmfamily]{Ligatures=TeX,Scale=1}
\fi
\usepackage{lmodern}
\ifPDFTeX\else
  % xetex/luatex font selection
\fi
% Use upquote if available, for straight quotes in verbatim environments
\IfFileExists{upquote.sty}{\usepackage{upquote}}{}
\IfFileExists{microtype.sty}{% use microtype if available
  \usepackage[]{microtype}
  \UseMicrotypeSet[protrusion]{basicmath} % disable protrusion for tt fonts
}{}
\makeatletter
\@ifundefined{KOMAClassName}{% if non-KOMA class
  \IfFileExists{parskip.sty}{%
    \usepackage{parskip}
  }{% else
    \setlength{\parindent}{0pt}
    \setlength{\parskip}{6pt plus 2pt minus 1pt}}
}{% if KOMA class
  \KOMAoptions{parskip=half}}
\makeatother
\usepackage{listings}
\newcommand{\passthrough}[1]{#1}
\lstset{defaultdialect=[5.3]Lua}
\lstset{defaultdialect=[x86masm]Assembler}
\usepackage{longtable,booktabs,array}
\newcounter{none} % for unnumbered tables
\usepackage{calc} % for calculating minipage widths
% Correct order of tables after \paragraph or \subparagraph
\usepackage{etoolbox}
\makeatletter
\patchcmd\longtable{\par}{\if@noskipsec\mbox{}\fi\par}{}{}
\makeatother
% Allow footnotes in longtable head/foot
\IfFileExists{footnotehyper.sty}{\usepackage{footnotehyper}}{\usepackage{footnote}}
\makesavenoteenv{longtable}
\setlength{\emergencystretch}{3em} % prevent overfull lines
\providecommand{\tightlist}{%
  \setlength{\itemsep}{0pt}\setlength{\parskip}{0pt}}
\usepackage{bookmark}
\IfFileExists{xurl.sty}{\usepackage{xurl}}{} % add URL line breaks if available
\urlstyle{same}
% fallback for those not using the hyperref driver hyperxmp:
\makeatletter
\@ifundefined{xmpquote}{\newcommand{\xmpquote}[1]{#1}}{}
\makeatother
\hypersetup{
  colorlinks=true,
  linkcolor={blue},
  filecolor={Maroon},
  citecolor={Blue},
  urlcolor={blue},
  pdfcreator={LaTeX via pandoc}}

\author{}
\date{}

\begin{document}

\section{计算机体系结构期末复习宝典}\label{ux8ba1ux7b97ux673aux4f53ux7cfbux7ed3ux6784ux671fux672bux590dux4e60ux5b9dux5178}

整理依据：当前文件夹内第 0-8 章课件，以及三次作业 PDF 的课程范围。作业
PDF
抽取文本为空，疑似扫描/图片型，因此例题按课件中的核心公式、典型课堂例题和常见考试题型重构。

\subsection{0. 考前抓大放小}\label{ux8003ux524dux6293ux5927ux653eux5c0f}

最容易考计算题的部分：

\begin{enumerate}
\def\labelenumi{\arabic{enumi}.}
\tightlist
\item
  性能评价：CPU time、CPI、MIPS、Amdahl 定律。
\item
  指令系统：栈/累加器/寄存器-存储器/Load-Store
  结构的代码序列比较，寻址方式，MIPS 指令格式。
\item
  流水线：加速比、吞吐率、效率、结构/数据/控制冲突，RAW/WAR/WAW，旁路与停顿。
\item
  指令级并行：静态调度、循环展开、记分牌、Tomasulo、寄存器重命名。
\item
  Cache：地址划分、映像方式、失效率、AMAT、多级 Cache、3C 失效。
\item
  I/O 与可靠性：Amdahl 在 I/O 中的应用，MTTF/MTTR/Availability，RAID
  级别。
\end{enumerate}

最容易考概念辨析的部分：

\begin{enumerate}
\def\labelenumi{\arabic{enumi}.}
\tightlist
\item
  体系结构、组成、实现三者区别。
\item
  RISC/CISC，Load-Store 结构优缺点。
\item
  Flynn 分类：SISD、SIMD、MISD、MIMD。
\item
  Cache 写策略：写直达/写回，写分配/不写分配。
\item
  RAID 0/1/3/4/5/6/10 的性能、容量和容错。
\end{enumerate}

\subsection{1.
计算机系统结构基础与性能评价}\label{ux8ba1ux7b97ux673aux7cfbux7edfux7ed3ux6784ux57faux7840ux4e0eux6027ux80fdux8bc4ux4ef7}

\subsubsection{1.1 必背知识点}\label{ux5fc5ux80ccux77e5ux8bc6ux70b9}

计算机系统结构是软硬件之间的接口，传统上常指
ISA，即机器语言程序员可见的属性，包括寄存器、寻址方式、指令格式、操作类型、异常和中断等。

体系结构、组成、实现的区别：

{\def\LTcaptype{none} % do not increment counter
\begin{longtable}[]{@{}
  >{\raggedright\arraybackslash}p{(\linewidth - 4\tabcolsep) * \real{0.3333}}
  >{\raggedright\arraybackslash}p{(\linewidth - 4\tabcolsep) * \real{0.3333}}
  >{\raggedright\arraybackslash}p{(\linewidth - 4\tabcolsep) * \real{0.3333}}@{}}
\toprule\noalign{}
\begin{minipage}[b]{\linewidth}\raggedright
层次
\end{minipage} & \begin{minipage}[b]{\linewidth}\raggedright
关注点
\end{minipage} & \begin{minipage}[b]{\linewidth}\raggedright
例子
\end{minipage} \\
\midrule\noalign{}
\endhead
\bottomrule\noalign{}
\endlastfoot
体系结构 Architecture & 程序员可见功能 & 是否有乘法指令 \\
组成 Organization & 如何用逻辑结构实现 ISA &
用专用乘法器还是加法器迭代实现乘法 \\
实现 Implementation & 物理实现 & 器件、工艺、封装、时钟、电源、散热 \\
\end{longtable}
}

性能基本公式：

\begin{lstlisting}
CPU time = Instruction Count × CPI × Clock Cycle Time
CPU time = Instruction Count × CPI / Clock Rate
CPI = Σ(指令比例_i × CPI_i)
Speedup = 原执行时间 / 新执行时间
\end{lstlisting}

Amdahl 定律：

\begin{lstlisting}
Speedup = 1 / ((1 - f) + f / S)
\end{lstlisting}

其中 \passthrough{\lstinline!f!}
是可改进部分在原执行时间中的比例，\passthrough{\lstinline!S!}
是该部分加速倍数。极限加速比为 \passthrough{\lstinline!1 / (1 - f)!}。

Flynn 分类：

{\def\LTcaptype{none} % do not increment counter
\begin{longtable}[]{@{}lll@{}}
\toprule\noalign{}
类型 & 含义 & 例子 \\
\midrule\noalign{}
\endhead
\bottomrule\noalign{}
\endlastfoot
SISD & 单指令流单数据流 & 传统单核标量机 \\
SIMD & 单指令流多数据流 & 向量机、GPU 的部分执行模型 \\
MISD & 多指令流单数据流 & 很少见 \\
MIMD & 多指令流多数据流 & 多核、多处理器 \\
\end{longtable}
}

\subsubsection{1.2 例题：Amdahl
定律}\label{ux4f8bux9898amdahl-ux5b9aux5f8b}

题目：某程序中浮点运算占原执行时间的 40\%。若把浮点部件加速 20
倍，程序总加速比是多少？如果浮点部件无限快，最大加速比是多少？

讲解：

\begin{lstlisting}
f = 0.4, S = 20
Speedup = 1 / ((1 - 0.4) + 0.4 / 20)
        = 1 / (0.6 + 0.02)
        = 1 / 0.62
        = 1.613
\end{lstlisting}

若浮点部件无限快，\passthrough{\lstinline!S -> ∞!}：

\begin{lstlisting}
Speedup_max = 1 / (1 - 0.4) = 1 / 0.6 = 1.667
\end{lstlisting}

结论：局部优化会受到未优化部分限制。考试中看到``某部分占比 +
加速倍数''，优先套 Amdahl。

\subsubsection{1.3 例题：CPU
性能比较}\label{ux4f8bux9898cpu-ux6027ux80fdux6bd4ux8f83}

题目：机器 A 的时钟周期为 1 ns，CPI 为 2.0；机器 B 的时钟周期为 1.5
ns，CPI 为 1.2。两者执行同一程序，指令条数相同，谁更快？

讲解：

只比较 \passthrough{\lstinline!CPI × Clock Cycle Time!}：

\begin{lstlisting}
A: 2.0 × 1 ns = 2.0 ns/指令
B: 1.2 × 1.5 ns = 1.8 ns/指令
\end{lstlisting}

B 更快，加速比：

\begin{lstlisting}
Speedup_B_over_A = 2.0 / 1.8 = 1.11
\end{lstlisting}

注意：主频高不一定快，必须同时看 CPI 和指令条数。

\subsection{2. 指令系统结构
ISA}\label{ux6307ux4ee4ux7cfbux7edfux7ed3ux6784-isa}

\subsubsection{2.1 必背知识点}\label{ux5fc5ux80ccux77e5ux8bc6ux70b9-1}

指令系统常见操作数组织方式：

{\def\LTcaptype{none} % do not increment counter
\begin{longtable}[]{@{}
  >{\raggedright\arraybackslash}p{(\linewidth - 4\tabcolsep) * \real{0.3333}}
  >{\raggedright\arraybackslash}p{(\linewidth - 4\tabcolsep) * \real{0.3333}}
  >{\raggedright\arraybackslash}p{(\linewidth - 4\tabcolsep) * \real{0.3333}}@{}}
\toprule\noalign{}
\begin{minipage}[b]{\linewidth}\raggedright
结构
\end{minipage} & \begin{minipage}[b]{\linewidth}\raggedright
特点
\end{minipage} & \begin{minipage}[b]{\linewidth}\raggedright
\passthrough{\lstinline!Z = X + Y!} 示例
\end{minipage} \\
\midrule\noalign{}
\endhead
\bottomrule\noalign{}
\endlastfoot
栈结构 & 操作数隐含在栈顶 &
\passthrough{\lstinline!push X; push Y; add; pop Z!} \\
累加器结构 & 一个隐含累加器 &
\passthrough{\lstinline!load X; add Y; store Z!} \\
寄存器-存储器 RM & ALU 可直接访问存储器 &
\passthrough{\lstinline!load R1,X; add R1,Y; store R1,Z!} \\
Load-Store/RR & ALU 只操作寄存器 &
\passthrough{\lstinline!load R1,X; load R2,Y; add R3,R1,R2; store R3,Z!} \\
\end{longtable}
}

RISC 常见特征：

\begin{enumerate}
\def\labelenumi{\arabic{enumi}.}
\tightlist
\item
  指令格式规整，长度固定。
\item
  Load-Store 结构，访存指令与 ALU 指令分离。
\item
  大量通用寄存器。
\item
  寻址方式较少。
\item
  便于流水线实现。
\end{enumerate}

MIPS 常见格式：

{\def\LTcaptype{none} % do not increment counter
\begin{longtable}[]{@{}lll@{}}
\toprule\noalign{}
格式 & 用途 & 关键字段 \\
\midrule\noalign{}
\endhead
\bottomrule\noalign{}
\endlastfoot
R 型 & 寄存器-寄存器运算 &
\passthrough{\lstinline!op, rs, rt, rd, shamt, funct!} \\
I 型 & 立即数、访存、分支 &
\passthrough{\lstinline!op, rs, rt, immediate!} \\
J 型 & 跳转 & \passthrough{\lstinline!op, target!} \\
\end{longtable}
}

常见寻址方式：

\begin{enumerate}
\def\labelenumi{\arabic{enumi}.}
\tightlist
\item
  立即寻址：操作数在指令中。
\item
  寄存器寻址：操作数在寄存器中。
\item
  基址/偏移寻址：\passthrough{\lstinline!Mem[Register + offset]!}，MIPS
  访存常用。
\item
  PC 相对寻址：分支目标相对 PC。
\item
  伪直接寻址：跳转指令常见。
\end{enumerate}

\subsubsection{2.2 例题：不同 ISA
写代码}\label{ux4f8bux9898ux4e0dux540c-isa-ux5199ux4ee3ux7801}

题目：设 \passthrough{\lstinline!X, Y, Z!} 均在存储器中，写出
\passthrough{\lstinline!Z = X + Y!}
在栈、累加器、寄存器-存储器、Load-Store 四类机器上的代码。

讲解：

\begin{lstlisting}
# 栈结构
push X
push Y
add
pop Z

# 累加器结构
load X
add Y
store Z

# 寄存器-存储器结构
load R1, X
add  R1, Y
store R1, Z

# Load-Store 结构
load R1, X
load R2, Y
add  R3, R1, R2
store R3, Z
\end{lstlisting}

答题关键：Load-Store 结构的 ALU 指令不能直接用内存操作数，所以必须先
\passthrough{\lstinline!load!} 到寄存器。

\subsubsection{2.3 例题：平均 CPI}\label{ux4f8bux9898ux5e73ux5747-cpi}

题目：某程序指令构成为：ALU 50\%，Load 25\%，Store 10\%，Branch
15\%。各类 CPI 分别为 1、2、2、3。求平均 CPI。

讲解：

\begin{lstlisting}
CPI = 0.50×1 + 0.25×2 + 0.10×2 + 0.15×3
    = 0.50 + 0.50 + 0.20 + 0.45
    = 1.65
\end{lstlisting}

考试中比例必须先化为小数，最后单位是``周期/指令''。

\subsection{3. 流水线}\label{ux6d41ux6c34ux7ebf}

\subsubsection{3.1 必背知识点}\label{ux5fc5ux80ccux77e5ux8bc6ux70b9-2}

流水线的目标是提高吞吐率，而不是降低单条指令的延迟。理想
\passthrough{\lstinline!k!} 级流水线执行 \passthrough{\lstinline!n!}
条指令：

\begin{lstlisting}
非流水时间 = n × k × Δt
流水时间 = (k + n - 1) × Δt
加速比 = n×k / (k + n - 1)
n 很大时，加速比趋近 k
吞吐率 = n / ((k+n-1)Δt)
效率 = 加速比 / k
\end{lstlisting}

流水线冲突：

{\def\LTcaptype{none} % do not increment counter
\begin{longtable}[]{@{}lll@{}}
\toprule\noalign{}
冲突 & 原因 & 解决 \\
\midrule\noalign{}
\endhead
\bottomrule\noalign{}
\endlastfoot
结构冲突 & 资源不够 & 增加资源、分时、停顿 \\
数据冲突 RAW & 后指令读前指令结果 & 旁路/转发、停顿、调度 \\
数据冲突 WAR & 后指令先写，破坏前指令读 & 寄存器重命名、顺序读 \\
数据冲突 WAW & 后指令先写，破坏最终写序 & 寄存器重命名、顺序提交 \\
控制冲突 & 分支改变 PC & 预测、延迟槽、提前判断 \\
\end{longtable}
}

五级 MIPS 流水线：

\begin{lstlisting}
IF 取指 -> ID 译码/读寄存器 -> EX 执行/地址计算 -> MEM 访存 -> WB 写回
\end{lstlisting}

\subsubsection{3.2
例题：流水线加速比}\label{ux4f8bux9898ux6d41ux6c34ux7ebfux52a0ux901fux6bd4}

题目：一条指令分为 5 个阶段，每阶段 1 ns。执行 100
条指令，理想流水线相对非流水线的加速比是多少？

讲解：

\begin{lstlisting}
非流水时间 = 100 × 5 × 1 = 500 ns
流水时间 = (5 + 100 - 1) × 1 = 104 ns
加速比 = 500 / 104 = 4.81
\end{lstlisting}

如果指令数趋于无穷，加速比才接近 5。不要直接写 5。

\subsubsection{3.3
例题：数据相关判断}\label{ux4f8bux9898ux6570ux636eux76f8ux5173ux5224ux65ad}

题目：判断以下指令间的数据相关类型。

\begin{lstlisting}
I1: ADD R1, R2, R3
I2: SUB R4, R1, R5
I3: OR  R1, R6, R7
I4: AND R8, R1, R9
\end{lstlisting}

讲解：

\begin{enumerate}
\def\labelenumi{\arabic{enumi}.}
\tightlist
\item
  \passthrough{\lstinline!I1 -> I2!}：\passthrough{\lstinline!I1!} 写
  R1，\passthrough{\lstinline!I2!} 读 R1，是 RAW 真相关。
\item
  \passthrough{\lstinline!I2 -> I3!}：无同一寄存器读写相关。
\item
  \passthrough{\lstinline!I1 -> I3!}：两者都写 R1，是 WAW 输出相关。
\item
  \passthrough{\lstinline!I2 -> I3!}：\passthrough{\lstinline!I2!} 读
  R1，\passthrough{\lstinline!I3!} 写 R1，是 WAR 反相关。
\item
  \passthrough{\lstinline!I3 -> I4!}：\passthrough{\lstinline!I3!} 写
  R1，\passthrough{\lstinline!I4!} 读 R1，是 RAW。
\end{enumerate}

解题口诀：先看同一个寄存器，再看``前后指令''的读写顺序：写后读
RAW，读后写 WAR，写后写 WAW。

\subsubsection{3.4 例题：分支惩罚下的
CPI}\label{ux4f8bux9898ux5206ux652fux60e9ux7f5aux4e0bux7684-cpi}

题目：理想流水线 CPI 为 1，分支指令占 20\%，分支预测失败率
10\%，预测失败惩罚 3 个周期。求实际 CPI。

讲解：

\begin{lstlisting}
额外 CPI = 分支比例 × 失败率 × 惩罚
         = 0.20 × 0.10 × 3
         = 0.06
实际 CPI = 1 + 0.06 = 1.06
\end{lstlisting}

控制冲突题目通常就是``发生频率 × 代价''。

\subsection{4. 向量处理机}\label{ux5411ux91cfux5904ux7406ux673a}

\subsubsection{4.1 必背知识点}\label{ux5fc5ux80ccux77e5ux8bc6ux70b9-3}

向量处理适合对数组、大规模数据做相同操作，核心是开发数据级并行。

常见结构：

\begin{enumerate}
\def\labelenumi{\arabic{enumi}.}
\tightlist
\item
  存储器-存储器型向量机：向量操作数直接来自存储器，结果也回存储器。
\item
  寄存器-寄存器型向量机：向量先装入向量寄存器，运算后再写回存储器，典型如
  CRAY-1。
\end{enumerate}

向量执行常见概念：

\begin{enumerate}
\def\labelenumi{\arabic{enumi}.}
\tightlist
\item
  向量长度 VL：一次向量指令处理的元素个数。
\item
  启动时间 startup time：流水线装满前的开销。
\item
  链接
  chaining：前一功能部件产生的元素可直接送入后一功能部件，减少等待。
\item
  分段开采 strip mining：当数组长度大于最大向量长度时，分多段处理。
\end{enumerate}

\subsubsection{4.2
例题：分段开采}\label{ux4f8bux9898ux5206ux6bb5ux5f00ux91c7}

题目：某向量寄存器最大长度 MVL = 64，要处理长度 N = 200
的数组。需要几段？每段 VL 是多少？

讲解：

\begin{lstlisting}
200 = 64 + 64 + 64 + 8
\end{lstlisting}

需要 4 段，前三段 \passthrough{\lstinline!VL = 64!}，最后一段
\passthrough{\lstinline!VL = 8!}。

答题时说明：循环每次处理
\passthrough{\lstinline!min(MVL, 剩余元素数)!}。

\subsubsection{4.3
例题：向量机为什么快}\label{ux4f8bux9898ux5411ux91cfux673aux4e3aux4ec0ux4e48ux5feb}

题目：为什么向量处理机比标量循环更适合
\passthrough{\lstinline!A[i] = B[i] + C[i]!}？

讲解：

标量循环每次迭代都需要取指、译码、分支判断和循环控制。向量指令把多个元素的相同操作合成一条指令，减少指令条数和分支开销，并可让功能部件流水化。若有多个存储体，还可连续供数，提高吞吐率。

\subsection{5. 指令级并行、动态调度与
Tomasulo}\label{ux6307ux4ee4ux7ea7ux5e76ux884cux52a8ux6001ux8c03ux5ea6ux4e0e-tomasulo}

\subsubsection{5.1 必背知识点}\label{ux5fc5ux80ccux77e5ux8bc6ux70b9-4}

指令级并行 ILP：

\begin{lstlisting}
ILP = 指令数 / 所需周期数
IPC = Instructions Per Cycle
CPI = 1 / IPC
\end{lstlisting}

流水线实际 CPI：

\begin{lstlisting}
CPI = 理想 CPI + 结构冲突停顿 + 数据冲突停顿 + 控制冲突停顿
\end{lstlisting}

静态调度由编译器完成，常见方法：

\begin{enumerate}
\def\labelenumi{\arabic{enumi}.}
\tightlist
\item
  指令重排：把无关指令插入相关指令之间。
\item
  循环展开：减少分支开销，暴露更多独立指令。
\item
  寄存器重命名：消除 WAR/WAW 名相关。
\end{enumerate}

动态调度由硬件完成：

{\def\LTcaptype{none} % do not increment counter
\begin{longtable}[]{@{}
  >{\raggedright\arraybackslash}p{(\linewidth - 4\tabcolsep) * \real{0.3333}}
  >{\raggedright\arraybackslash}p{(\linewidth - 4\tabcolsep) * \real{0.3333}}
  >{\raggedright\arraybackslash}p{(\linewidth - 4\tabcolsep) * \real{0.3333}}@{}}
\toprule\noalign{}
\begin{minipage}[b]{\linewidth}\raggedright
方法
\end{minipage} & \begin{minipage}[b]{\linewidth}\raggedright
核心思想
\end{minipage} & \begin{minipage}[b]{\linewidth}\raggedright
解决问题
\end{minipage} \\
\midrule\noalign{}
\endhead
\bottomrule\noalign{}
\endlastfoot
记分牌 Scoreboard & 集中记录功能部件和寄存器状态 & 允许乱序执行，但遇
WAR/WAW 可能停顿 \\
Tomasulo & 保留站 + 公共数据总线 CDB + 寄存器重命名 & 消除 WAR/WAW，减少
RAW 等待 \\
ROB & 重排序缓冲 & 支持精确异常和按程序顺序提交 \\
\end{longtable}
}

Tomasulo 三个阶段：

\begin{enumerate}
\def\labelenumi{\arabic{enumi}.}
\tightlist
\item
  Issue：若保留站空闲，发射；源操作数可用则读值，否则记录生产者 tag。
\item
  Execute：源操作数都就绪后执行。
\item
  Write result：通过 CDB 广播结果，等待者按 tag 捕获。
\end{enumerate}

\subsubsection{5.2
例题：循环调度}\label{ux4f8bux9898ux5faaux73afux8c03ux5ea6}

题目：有如下循环，假设 \passthrough{\lstinline!L.D -> ADD.D!} 需要 1
个停顿，\passthrough{\lstinline!ADD.D -> S.D!} 需要 2 个停顿，分支也有 1
个停顿。对代码做简单调度。

\begin{lstlisting}
Loop:
L.D    F0, 0(R1)
ADD.D  F4, F0, F2
S.D    F4, 0(R1)
DADDIU R1, R1, #-8
BNE    R1, R2, Loop
\end{lstlisting}

讲解：

可把与浮点数据无关的地址更新和分支提前，填补部分空隙：

\begin{lstlisting}
Loop:
L.D    F0, 0(R1)
DADDIU R1, R1, #-8
ADD.D  F4, F0, F2
BNE    R1, R2, Loop
S.D    F4, 8(R1)
\end{lstlisting}

为什么 \passthrough{\lstinline!S.D!} 地址变成
\passthrough{\lstinline!8(R1)!}：因为 \passthrough{\lstinline!R1!}
已经提前减了 8，要存回原来的地址，必须用 \passthrough{\lstinline!8(R1)!}
修正。

\subsubsection{5.3
例题：循环展开}\label{ux4f8bux9898ux5faaux73afux5c55ux5f00}

题目：把上题循环展开 4 次时，为什么要使用
\passthrough{\lstinline!F0/F4, F6/F8, F10/F12, F14/F16!} 等不同寄存器？

讲解：

如果四次展开都使用同一组寄存器，会产生大量 WAW 和 WAR
名相关，限制重排。使用不同寄存器相当于手工寄存器重命名，可以让四次迭代的
Load、Add、Store 互相独立，编译器能把多个 Load 放在前面，再集中执行 Add
和 Store，从而隐藏延迟。

\subsubsection{5.4 例题：Tomasulo
消除名相关}\label{ux4f8bux9898tomasulo-ux6d88ux9664ux540dux76f8ux5173}

题目：判断 Tomasulo 如何处理下面的相关：

\begin{lstlisting}
I1: DIV.D F6, F0, F2
I2: ADD.D F8, F6, F4
I3: SUB.D F6, F10, F12
\end{lstlisting}

讲解：

\begin{enumerate}
\def\labelenumi{\arabic{enumi}.}
\tightlist
\item
  \passthrough{\lstinline!I1 -> I2!}：\passthrough{\lstinline!I2!} 读
  \passthrough{\lstinline!I1!} 产生的 F6，是 RAW
  真相关，不能消除，必须等 \passthrough{\lstinline!I1!} 结果。
\item
  \passthrough{\lstinline!I1 -> I3!}：两者都写 F6，是 WAW
  名相关。Tomasulo 给两个写 F6 的指令分配不同保留站
  tag，寄存器结果状态只保留最新生产者，因此可避免错误覆盖。
\item
  \passthrough{\lstinline!I2 -> I3!}：\passthrough{\lstinline!I2!} 读
  F6，\passthrough{\lstinline!I3!} 写 F6，是 WAR 名相关。由于
  \passthrough{\lstinline!I2!} 需要的是 \passthrough{\lstinline!I1!} 的
  tag，而不是``寄存器 F6 这个名字''，所以 \passthrough{\lstinline!I3!}
  改写 F6 的名字不会破坏 \passthrough{\lstinline!I2!}。
\end{enumerate}

结论：Tomasulo 不能消除 RAW 真相关，但能通过寄存器重命名消除 WAR/WAW。

\subsection{6. Cache
与存储层次}\label{cache-ux4e0eux5b58ux50a8ux5c42ux6b21}

\subsubsection{6.1 必背知识点}\label{ux5fc5ux80ccux77e5ux8bc6ux70b9-5}

存储层次建立在局部性原理上：

\begin{enumerate}
\def\labelenumi{\arabic{enumi}.}
\tightlist
\item
  时间局部性：最近访问过的数据/指令很可能再次访问。
\item
  空间局部性：访问某地址后，邻近地址很可能被访问。
\end{enumerate}

Cache 基本问题：

{\def\LTcaptype{none} % do not increment counter
\begin{longtable}[]{@{}ll@{}}
\toprule\noalign{}
问题 & 选项 \\
\midrule\noalign{}
\endhead
\bottomrule\noalign{}
\endlastfoot
放哪里 & 直接映像、组相联、全相联 \\
找不找得到 & tag + valid 位比较 \\
满了替换谁 & 随机、FIFO、LRU \\
写命中怎么办 & 写直达、写回 \\
写不命中怎么办 & 写分配、不写分配 \\
\end{longtable}
}

映像方式：

\begin{lstlisting}
直接映像：cache line = memory block mod cache line 数
组相联：set = memory block mod set 数
全相联：任意块可放任意行
\end{lstlisting}

地址划分：

\begin{lstlisting}
块内偏移位 = log2(块大小)
索引位 = log2(Cache 行数或组数)
Tag 位 = 地址总位数 - 索引位 - 块内偏移位
\end{lstlisting}

平均访存时间 AMAT：

\begin{lstlisting}
AMAT = Hit Time + Miss Rate × Miss Penalty
\end{lstlisting}

多级 Cache：

\begin{lstlisting}
AMAT = L1 Hit Time + L1 Miss Rate × (L2 Hit Time + L2 Miss Rate × Memory Penalty)
\end{lstlisting}

注意：二级 Cache 的 miss rate
可能是局部失效率，也可能是全局失效率。全局失效率 = L1 miss rate × L2
local miss rate。

3C 失效：

{\def\LTcaptype{none} % do not increment counter
\begin{longtable}[]{@{}lll@{}}
\toprule\noalign{}
类型 & 原因 & 降低方法 \\
\midrule\noalign{}
\endhead
\bottomrule\noalign{}
\endlastfoot
Compulsory 强制失效 & 第一次访问 & 预取、增大块 \\
Capacity 容量失效 & Cache 装不下工作集 & 增大 Cache \\
Conflict 冲突失效 & 映像位置冲突 & 提高相联度、Victim Cache \\
\end{longtable}
}

\subsubsection{6.2 例题：Cache
地址划分}\label{ux4f8bux9898cache-ux5730ux5740ux5212ux5206}

题目：32 位地址，Cache 容量 16 KB，块大小 32 B，直接映像。求
offset、index、tag 位数。

讲解：

\begin{lstlisting}
块大小 = 32 B = 2^5 B，所以 offset = 5 位
Cache 行数 = 16 KB / 32 B = 16384 / 32 = 512 = 2^9，所以 index = 9 位
tag = 32 - 9 - 5 = 18 位
\end{lstlisting}

答案：\passthrough{\lstinline!tag 18 位，index 9 位，offset 5 位!}。

\subsubsection{6.3
例题：组相联地址划分}\label{ux4f8bux9898ux7ec4ux76f8ux8054ux5730ux5740ux5212ux5206}

题目：32 位地址，Cache 容量 64 KB，块大小 64 B，4 路组相联。求组索引位和
tag 位。

讲解：

\begin{lstlisting}
offset = log2(64) = 6 位
Cache 总行数 = 64 KB / 64 B = 1024 行
组数 = 1024 / 4 = 256 = 2^8
index = 8 位
tag = 32 - 8 - 6 = 18 位
\end{lstlisting}

注意：组相联的 index 看``组数''，不是总行数。

\subsubsection{6.4 例题：AMAT}\label{ux4f8bux9898amat}

题目：L1 命中时间 1 cycle，失效率 5\%，失效代价 50 cycles。求 AMAT。

讲解：

\begin{lstlisting}
AMAT = 1 + 0.05 × 50 = 3.5 cycles
\end{lstlisting}

若题目问``每条指令平均访存 1.3 次''，则内存停顿 CPI 为：

\begin{lstlisting}
Memory stall CPI = 1.3 × 0.05 × 50 = 3.25
\end{lstlisting}

看清题目是``每次访存''还是``每条指令''。

\subsubsection{6.5 例题：多级
Cache}\label{ux4f8bux9898ux591aux7ea7-cache}

题目：L1 命中时间 1 cycle，L1 失效率 4\%；L2 命中时间 8 cycles，L2
局部失效率 20\%；主存代价 100 cycles。求 AMAT。

讲解：

\begin{lstlisting}
L2 平均服务时间 = 8 + 0.20 × 100 = 28 cycles
AMAT = 1 + 0.04 × 28 = 2.12 cycles
\end{lstlisting}

如果题目给的是 L2 全局失效率 0.8\%，则公式应写成：

\begin{lstlisting}
AMAT = L1 hit time + L1 miss rate × L2 hit time + L2 global miss rate × memory penalty
\end{lstlisting}

\subsubsection{6.6
例题：直接映像冲突}\label{ux4f8bux9898ux76f4ux63a5ux6620ux50cfux51b2ux7a81}

题目：Cache 有 8 行，块大小 1 word，访问块序列
\passthrough{\lstinline!0, 8, 0, 8, 0, 8!}，直接映像 Cache
初始为空，命中率是多少？

讲解：

直接映像行号：

\begin{lstlisting}
line = block number mod 8
0 mod 8 = 0
8 mod 8 = 0
\end{lstlisting}

块 0 和块 8 总是争同一行。访问过程全是失效：

\begin{lstlisting}
0 miss, 8 miss, 0 miss, 8 miss, 0 miss, 8 miss
\end{lstlisting}

命中率 \passthrough{\lstinline!0/6 = 0\%!}。这是典型 conflict miss。

\subsection{7. I/O、可靠性与
RAID}\label{ioux53efux9760ux6027ux4e0e-raid}

\subsubsection{7.1 必背知识点}\label{ux5fc5ux80ccux77e5ux8bc6ux70b9-6}

I/O 系统组成：

\begin{enumerate}
\def\labelenumi{\arabic{enumi}.}
\tightlist
\item
  I/O 设备：磁盘、SSD、键盘、网卡等。
\item
  I/O 接口/控制器：连接设备与系统总线。
\item
  I/O 软件：驱动、中断处理、操作系统 I/O 层。
\item
  DMA：设备与内存直接传输，减少 CPU 参与。
\end{enumerate}

I/O 性能常用指标：

\begin{enumerate}
\def\labelenumi{\arabic{enumi}.}
\tightlist
\item
  响应时间：单个请求从发出到完成的时间。
\item
  吞吐率：单位时间完成的 I/O 请求数或传输数据量。
\item
  利用率：设备忙碌时间比例。
\end{enumerate}

可靠性公式：

\begin{lstlisting}
MTTF = Mean Time To Failure，平均无故障时间
MTTR = Mean Time To Repair，平均修复时间
MTBF = MTTF + MTTR
Availability = MTTF / (MTTF + MTTR)
串联系统失效率 = 各部件失效率之和 = Σ(1/MTTF_i)
系统 MTTF = 1 / 系统失效率
\end{lstlisting}

RAID 速记：

{\def\LTcaptype{none} % do not increment counter
\begin{longtable}[]{@{}
  >{\raggedright\arraybackslash}p{(\linewidth - 10\tabcolsep) * \real{0.1667}}
  >{\raggedright\arraybackslash}p{(\linewidth - 10\tabcolsep) * \real{0.1667}}
  >{\raggedright\arraybackslash}p{(\linewidth - 10\tabcolsep) * \real{0.1667}}
  >{\raggedright\arraybackslash}p{(\linewidth - 10\tabcolsep) * \real{0.1667}}
  >{\raggedright\arraybackslash}p{(\linewidth - 10\tabcolsep) * \real{0.1667}}
  >{\raggedright\arraybackslash}p{(\linewidth - 10\tabcolsep) * \real{0.1667}}@{}}
\toprule\noalign{}
\begin{minipage}[b]{\linewidth}\raggedright
级别
\end{minipage} & \begin{minipage}[b]{\linewidth}\raggedright
核心
\end{minipage} & \begin{minipage}[b]{\linewidth}\raggedright
容错
\end{minipage} & \begin{minipage}[b]{\linewidth}\raggedright
读性能
\end{minipage} & \begin{minipage}[b]{\linewidth}\raggedright
写性能
\end{minipage} & \begin{minipage}[b]{\linewidth}\raggedright
容量利用率
\end{minipage} \\
\midrule\noalign{}
\endhead
\bottomrule\noalign{}
\endlastfoot
RAID 0 & 条带化 & 无 & 高 & 高 & 100\% \\
RAID 1 & 镜像 & 可坏一个镜像 & 高 & 中 & 50\% \\
RAID 3 & 位/字节级条带 + 专用校验盘 & 可坏 1 块 & 顺序读好 & 小写差 &
\passthrough{\lstinline!(N-1)/N!} \\
RAID 4 & 块级条带 + 专用校验盘 & 可坏 1 块 & 随机读好 & 校验盘瓶颈 &
\passthrough{\lstinline!(N-1)/N!} \\
RAID 5 & 块级条带 + 分布式校验 & 可坏 1 块 & 好 & 小写有读改写 &
\passthrough{\lstinline!(N-1)/N!} \\
RAID 6 & 双分布式校验 & 可坏 2 块 & 好 & 写代价更大 &
\passthrough{\lstinline!(N-2)/N!} \\
RAID 10 & 先镜像再条带 & 每组可坏 1 块 & 高 & 高 & 50\% \\
\end{longtable}
}

\subsubsection{7.2 例题：I/O 的 Amdahl
效应}\label{ux4f8bux9898io-ux7684-amdahl-ux6548ux5e94}

题目：某系统原来 I/O 时间占总执行时间 10\%。若 CPU 速度提高 10 倍，而
I/O 不变，总执行时间变为原来的多少？总体加速比是多少？

讲解：

设原总时间为 1：

\begin{lstlisting}
CPU 部分 = 0.90
I/O 部分 = 0.10
新时间 = 0.90 / 10 + 0.10 = 0.09 + 0.10 = 0.19
总加速比 = 1 / 0.19 = 5.26
\end{lstlisting}

结论：CPU 越快，I/O 越容易成为瓶颈。

\subsubsection{7.3 例题：系统 MTTF}\label{ux4f8bux9898ux7cfbux7edf-mttf}

题目：一个系统由 10 块磁盘、1 个控制器、1 个电源、1 个风扇组成。磁盘
MTTF 均为 1,000,000 小时，控制器 MTTF 为 500,000 小时，电源和风扇 MTTF
均为 200,000 小时。任一部件失效都会导致系统失效，求系统 MTTF。

讲解：

串联系统失效率相加：

\begin{lstlisting}
λ = 10/1,000,000 + 1/500,000 + 1/200,000 + 1/200,000
  = 0.000010 + 0.000002 + 0.000005 + 0.000005
  = 0.000022
MTTF_system = 1 / 0.000022 = 45,454.5 小时
\end{lstlisting}

若题目还有 1 个 SCSI 线缆且 MTTF 为 1,000,000 小时，则再加
\passthrough{\lstinline!1/1,000,000!}，系统 MTTF 约为：

\begin{lstlisting}
λ = 0.000023
MTTF = 43,478 小时
\end{lstlisting}

\subsubsection{7.4 例题：可用性}\label{ux4f8bux9898ux53efux7528ux6027}

题目：某设备 MTTF = 1000 小时，MTTR = 2 小时，求 Availability。

讲解：

\begin{lstlisting}
Availability = MTTF / (MTTF + MTTR)
             = 1000 / 1002
             = 0.9980
             = 99.80%
\end{lstlisting}

可用性高不等于永不坏，也取决于修复时间。

\subsubsection{7.5 例题：RAID 容量}\label{ux4f8bux9898raid-ux5bb9ux91cf}

题目：8 块 4 TB 磁盘分别组成 RAID 0、RAID 1、RAID 5、RAID 6、RAID
10，有效容量是多少？

讲解：

\begin{lstlisting}
RAID 0: 8 × 4 = 32 TB
RAID 1: 需要镜像，容量约 32 / 2 = 16 TB
RAID 5: 损失 1 块校验盘容量，(8 - 1) × 4 = 28 TB
RAID 6: 损失 2 块校验盘容量，(8 - 2) × 4 = 24 TB
RAID 10: 镜像后条带，容量约 32 / 2 = 16 TB
\end{lstlisting}

RAID 0 容量和性能高，但没有容错；RAID 6 容错强但写代价更高。

\subsection{8. 作业专题强化}\label{ux4f5cux4e1aux4e13ux9898ux5f3aux5316}

这一章按三次作业来复盘。作业题比课件概念更接近考试风格，尤其是性能评价、流水线时空图、预约表、记分牌和
Tomasulo。

\subsubsection{8.1
第一次作业：性能评价与平均值}\label{ux7b2cux4e00ux6b21ux4f5cux4e1aux6027ux80fdux8bc4ux4ef7ux4e0eux5e73ux5747ux503c}

题目来源：教材第 1 章课后习题
\passthrough{\lstinline!1.2、1.7、1.8、1.10、1.11!}。

作业覆盖点：

\begin{enumerate}
\def\labelenumi{\arabic{enumi}.}
\tightlist
\item
  体系结构、计算机组成、计算机实现的区别。
\item
  指令条数、CPI、MIPS、程序执行时间。
\item
  算术平均、几何平均、调和平均的适用场景。
\item
  改进某类指令 CPI 后的总 CPI 与加速比。
\end{enumerate}

\paragraph{例题 1：CPI、MIPS
与执行时间}\label{ux4f8bux9898-1cpimips-ux4e0eux6267ux884cux65f6ux95f4}

题目：某程序有四类指令，条数分别为 45000、8000、75000、1500，各类 CPI
分别为 1、4、2、2。处理器主频 400 MHz。求总指令条数、平均 CPI、MIPS
和执行时间。

讲解：

\begin{lstlisting}
IC = 45000 + 8000 + 75000 + 1500 = 129500
总周期数 = 45000×1 + 8000×4 + 75000×2 + 1500×2
       = 230000
平均 CPI = 230000 / 129500 = 1.776
MIPS = 主频(MHz) / CPI = 400 / 1.776 = 225.2 MIPS
执行时间 = 总周期数 / 主频 = 230000 / (400×10^6)
       = 5.75×10^-4 s = 575 us
\end{lstlisting}

易错点：MIPS 是``每秒百万条指令''，但不能只看 MIPS
判断机器快慢，因为不同机器可能执行不同指令条数。

\paragraph{例题
2：三种平均值怎么选}\label{ux4f8bux9898-2ux4e09ux79cdux5e73ux5747ux503cux600eux4e48ux9009}

题目：三台机器 A、B、C 在四个程序上的 MIPS 如下：

{\def\LTcaptype{none} % do not increment counter
\begin{longtable}[]{@{}lrrr@{}}
\toprule\noalign{}
程序 & A & B & C \\
\midrule\noalign{}
\endhead
\bottomrule\noalign{}
\endlastfoot
1 & 100 & 10 & 5 \\
2 & 0.1 & 1 & 5 \\
3 & 0.2 & 0.1 & 2 \\
4 & 1 & 0.15 & 1 \\
\end{longtable}
}

讲解：

算术平均：

\begin{lstlisting}
A = (100 + 0.1 + 0.2 + 1) / 4 = 25.325
B = (10 + 1 + 0.1 + 0.15) / 4 = 2.8125
C = (5 + 5 + 2 + 1) / 4 = 3.25
\end{lstlisting}

几何平均：

\begin{lstlisting}
A = (100×0.1×0.2×1)^(1/4) = 1.189
B = (10×1×0.1×0.15)^(1/4) = 0.622
C = (5×5×2×1)^(1/4) = 2.659
\end{lstlisting}

调和平均：

\begin{lstlisting}
A = 4 / (1/100 + 1/0.1 + 1/0.2 + 1/1) = 0.250
B = 4 / (1/10 + 1/1 + 1/0.1 + 1/0.15) = 0.225
C = 4 / (1/5 + 1/5 + 1/2 + 1/1) = 2.105
\end{lstlisting}

考试判断：

\begin{enumerate}
\def\labelenumi{\arabic{enumi}.}
\tightlist
\item
  算术平均适合``同一机器运行等权程序时间''一类直接平均，但容易被极端值误导。
\item
  几何平均适合比较归一化后的相对性能。
\item
  调和平均适合平均``速率''，例如同一工作量下平均 MIPS。
\end{enumerate}

\paragraph{例题 3：改进某类指令
CPI}\label{ux4f8bux9898-3ux6539ux8fdbux67d0ux7c7bux6307ux4ee4-cpi}

题目：原程序中 30\% 指令 CPI 为 3，70\% 指令 CPI 为 1.25。若方案一把其中
40\% 指令的 CPI 降为 20/64，方案二把 70\% 指令 CPI 降为 1.5，哪个更好？

讲解：

原平均 CPI：

\begin{lstlisting}
CPI_old = 0.3×3 + 0.7×1.25 = 1.775
\end{lstlisting}

方案一：

\begin{lstlisting}
CPI_1 = 0.4×(20/64) + 0.6×3 = 1.925
\end{lstlisting}

方案二：

\begin{lstlisting}
CPI_2 = 0.3×3 + 0.7×1.5 = 1.95
\end{lstlisting}

若与某个原始基准 CPI 2.375 比较：

\begin{lstlisting}
Speedup_1 = 2.375 / 1.925 = 1.23
Speedup_2 = 2.375 / 1.95 = 1.22
\end{lstlisting}

答题关键：先写清楚``哪些指令比例被改进''，再算新 CPI，最后用
\passthrough{\lstinline!Speedup = CPI\_old / CPI\_new!}。

\subsubsection{8.2
第二次作业：流水线与预约表}\label{ux7b2cux4e8cux6b21ux4f5cux4e1aux6d41ux6c34ux7ebfux4e0eux9884ux7ea6ux8868}

题目来源：教材第 3 章课后习题
\passthrough{\lstinline!3.3、3.6、3.7、3.10、3.12!}。

作业覆盖点：

\begin{enumerate}
\def\labelenumi{\arabic{enumi}.}
\tightlist
\item
  分支处理：预测失败、预测成功、延迟分支。
\item
  非线性流水线预约表、禁止表、冲突向量。
\item
  吞吐率 TP、最大吞吐率 TPmax、效率 E。
\item
  功能部件重复设置对流水线性能的影响。
\item
  分支 CPI 对总加速比的影响。
\end{enumerate}

\paragraph{例题
1：分支处理三种办法}\label{ux4f8bux9898-1ux5206ux652fux5904ux7406ux4e09ux79cdux529eux6cd5}

题目：说明预测分支失败、预测分支成功、延迟分支三种方法。

讲解：

\begin{enumerate}
\def\labelenumi{\arabic{enumi}.}
\tightlist
\item
  预测分支失败：默认分支不发生，继续取顺序指令；如果实际发生分支，再清空错误取出的指令，转到目标地址。
\item
  预测分支成功：默认分支发生，优先从分支目标地址取指；如果实际不发生，再恢复顺序路径。
\item
  延迟分支：把分支后一条或几条指令作为延迟槽，无论分支是否发生都执行，编译器尽量填入有用指令。
\end{enumerate}

共同点：都在减少控制相关造成的流水线停顿。本质区别是``猜哪条路径''和``能否用编译器填槽''。

\paragraph{例题
2：线性流水线吞吐率与效率}\label{ux4f8bux9898-2ux7ebfux6027ux6d41ux6c34ux7ebfux541eux5410ux7387ux4e0eux6548ux7387}

题目：某 4 段流水线各段时间为 50 ns、50 ns、100 ns、200 ns，连续输入 10
个任务。求吞吐率和效率。

讲解：

流水线时钟周期由最慢段决定：

\begin{lstlisting}
Δt = max(50,50,100,200) = 200 ns
总时间 = (段数 + 任务数 - 1) × Δt = (4 + 10 - 1)×200 ns = 2600 ns
TP = 10 / 2600 ns = 1 / 260 ns
\end{lstlisting}

效率常写成：

\begin{lstlisting}
E = 实际忙碌时间 / 总可用流水段时间
  = 10×(50+50+100+200) / (4×2600)
  = 4000 / 10400 = 38.46%
\end{lstlisting}

如果采用细分瓶颈段或重复设置功能段，总时间会下降，效率也会变化。关键先画时空图。

\paragraph{例题
3：非线性流水线禁止表与冲突向量}\label{ux4f8bux9898-3ux975eux7ebfux6027ux6d41ux6c34ux7ebfux7981ux6b62ux8868ux4e0eux51b2ux7a81ux5411ux91cf}

题目：某非线性流水线预约表中，同一功能段被同一任务在时刻集合
\passthrough{\lstinline!\{1,3,6\}!} 使用，求这部分导致的禁止延迟。

讲解：

禁止延迟来自同一行中任意两个预约时刻的差：

\begin{lstlisting}
3 - 1 = 2
6 - 3 = 3
6 - 1 = 5
\end{lstlisting}

因此禁止表含
\passthrough{\lstinline!\{2,3,5\}!}。若多行预约表都给出时刻集合，则对每行都求差，再取并集。

冲突向量写法：

\begin{lstlisting}
C = (c_m ... c_2 c_1)
\end{lstlisting}

其中 \passthrough{\lstinline!c\_i = 1!} 表示延迟
\passthrough{\lstinline!i!} 禁止，\passthrough{\lstinline!c\_i = 0!}
表示允许。调度时每发射一个任务，就根据冲突向量决定下一次可用延迟。

\paragraph{例题 4：分支对 CPI
的三种情况}\label{ux4f8bux9898-4ux5206ux652fux5bf9-cpi-ux7684ux4e09ux79cdux60c5ux51b5}

题目：基础流水线 CPI 为 1，20\% 为条件分支，5\% 为无条件跳转。条件分支中
60\% 发生跳转。预测分支失败时发生分支需 2
周期惩罚，不发生无惩罚；无条件跳转惩罚 1 周期。求 CPI。

讲解：

\begin{lstlisting}
CPI = 1 + 条件分支比例×发生比例×惩罚 + 无条件跳转比例×惩罚
    = 1 + 0.20×0.60×2 + 0.05×1
    = 1.29
\end{lstlisting}

若预测分支成功，则条件分支不发生时才付惩罚：

\begin{lstlisting}
CPI = 1 + 0.20×0.40×2 + 0.05×1
    = 1.21
\end{lstlisting}

若延迟分支能填充一部分延迟槽，就把``未填满的比例 × 惩罚''作为额外 CPI。

\subsubsection{8.3 第三次作业：记分牌与
Tomasulo}\label{ux7b2cux4e09ux6b21ux4f5cux4e1aux8bb0ux5206ux724cux4e0e-tomasulo}

题目来源：第 3 次课后作业 PDF，题干直接给出记分牌与 Tomasulo 调度题。

作业覆盖点：

\begin{enumerate}
\def\labelenumi{\arabic{enumi}.}
\tightlist
\item
  记分牌算法思想。
\item
  Tomasulo 算法思想。
\item
  给定浮点指令序列，填写``流出/读操作数/执行/写回''周期。
\item
  在指定周期填写功能部件状态表和寄存器结果状态表。
\end{enumerate}

作业中使用的典型指令序列：

\begin{lstlisting}
L.D   F6, 34(R2)
L.D   F2, 45(R3)
ADD.D F8, F6, F2
MUL.D F4, F2, F6
SUB.D F8, F4, F2
ADD.D F4, F2, F6
\end{lstlisting}

假设功能部件：整数加法、浮点加法、浮点乘法、LOAD 部件各 2
个；LOAD、浮点加法、浮点乘法执行周期分别为 1、2、10，且尚未采用流水。

\paragraph{例题
1：记分牌算法怎么想}\label{ux4f8bux9898-1ux8bb0ux5206ux724cux7b97ux6cd5ux600eux4e48ux60f3}

记分牌核心思想：

\begin{enumerate}
\def\labelenumi{\arabic{enumi}.}
\tightlist
\item
  集中维护功能部件状态、寄存器结果状态和操作数就绪状态。
\item
  若没有结构冲突和 WAW，指令可以流出。
\item
  若源操作数可读且不会造成 WAR，进入读操作数。
\item
  操作数读完后执行。
\item
  写回时检查是否会破坏前面未读操作数的指令，避免 WAR。
\end{enumerate}

四阶段：

\begin{lstlisting}
Issue / 流出
Read Operands / 读操作数
Execute / 执行
Write Result / 写回
\end{lstlisting}

记分牌能乱序执行，但 WAR/WAW 仍可能造成停顿。

\paragraph{例题
2：记分牌时间表怎么看}\label{ux4f8bux9898-2ux8bb0ux5206ux724cux65f6ux95f4ux8868ux600eux4e48ux770b}

题目：对上述指令序列，作业中给出的记分牌调度结果为：

{\def\LTcaptype{none} % do not increment counter
\begin{longtable}[]{@{}lrrrr@{}}
\toprule\noalign{}
指令 & 流出 & 读操作数 & 执行完成 & 写回 \\
\midrule\noalign{}
\endhead
\bottomrule\noalign{}
\endlastfoot
\passthrough{\lstinline!L.D F6,34(R2)!} & 1 & 2 & 3 & 4 \\
\passthrough{\lstinline!L.D F2,45(R3)!} & 2 & 3 & 4 & 5 \\
\passthrough{\lstinline!ADD.D F8,F6,F2!} & 3 & 6 & 8 & 9 \\
\passthrough{\lstinline!MUL.D F4,F2,F6!} & 4 & 6 & 16 & 17 \\
\passthrough{\lstinline!SUB.D F8,F4,F2!} & 10 & 18 & 20 & 21 \\
\passthrough{\lstinline!ADD.D F4,F2,F6!} & 18 & 19 & 21 & 22 \\
\end{longtable}
}

讲解：

\begin{enumerate}
\def\labelenumi{\arabic{enumi}.}
\tightlist
\item
  两条 Load 最早流出，因为 Load 部件有 2 个，且目标寄存器不同。
\item
  \passthrough{\lstinline!ADD.D F8,F6,F2!} 要等 F6、F2
  都写回后才能读，因此读操作数在周期 6。
\item
  \passthrough{\lstinline!MUL.D F4,F2,F6!} 同样等 F2、F6，周期 6
  读，乘法执行 10 周期，所以周期 16 完成，17 写回。
\item
  \passthrough{\lstinline!SUB.D F8,F4,F2!} 写 F8，与前面
  \passthrough{\lstinline!ADD.D F8,...!} 有
  WAW，必须等前一条写回后才能流出，所以较晚。
\item
  \passthrough{\lstinline!ADD.D F4,F2,F6!} 写 F4，与
  \passthrough{\lstinline!MUL.D F4,...!} 有 WAW，也要等。
\end{enumerate}

考试中不要只背表，要会解释``为什么这条卡住''：结构冲突、RAW、WAR、WAW
分别找。

\paragraph{例题
3：记分牌状态表字段}\label{ux4f8bux9898-3ux8bb0ux5206ux724cux72b6ux6001ux8868ux5b57ux6bb5}

功能部件状态表常见字段：

{\def\LTcaptype{none} % do not increment counter
\begin{longtable}[]{@{}ll@{}}
\toprule\noalign{}
字段 & 含义 \\
\midrule\noalign{}
\endhead
\bottomrule\noalign{}
\endlastfoot
Busy & 功能部件是否忙 \\
Op & 当前操作 \\
Fi & 目的寄存器 \\
Fj, Fk & 源寄存器 \\
Qj, Qk & 产生源操作数的功能部件 \\
Rj, Rk & 源操作数是否已经就绪 \\
\end{longtable}
}

寄存器结果状态表：

\begin{lstlisting}
Result[Fi] = 将要写 Fi 的功能部件名
\end{lstlisting}

若某寄存器没有等待写回的生产者，则为空。

\paragraph{例题 4：Tomasulo
算法怎么想}\label{ux4f8bux9898-4tomasulo-ux7b97ux6cd5ux600eux4e48ux60f3}

Tomasulo 核心思想：

\begin{enumerate}
\def\labelenumi{\arabic{enumi}.}
\tightlist
\item
  用保留站保存等待执行的指令。
\item
  源操作数可用时记录值
  \passthrough{\lstinline!Vj/Vk!}；不可用时记录生产者 tag
  \passthrough{\lstinline!Qj/Qk!}。
\item
  CDB 广播结果，等待同一 tag 的保留站同时捕获。
\item
  寄存器结果状态保存 tag，从而实现寄存器重命名。
\end{enumerate}

所以 Tomasulo：

\begin{lstlisting}
RAW 真相关：不能消除，只能等值广播。
WAR/WAW 名相关：通过寄存器重命名消除。
\end{lstlisting}

作业中同一序列的 Tomasulo 调度结果为：

{\def\LTcaptype{none} % do not increment counter
\begin{longtable}[]{@{}lrrr@{}}
\toprule\noalign{}
指令 & 流出 & 执行完成 & 写回 \\
\midrule\noalign{}
\endhead
\bottomrule\noalign{}
\endlastfoot
\passthrough{\lstinline!L.D F6,34(R2)!} & 1 & 2 & 3 \\
\passthrough{\lstinline!L.D F2,45(R3)!} & 2 & 3 & 4 \\
\passthrough{\lstinline!ADD.D F8,F6,F2!} & 3 & 6 & 7 \\
\passthrough{\lstinline!MUL.D F4,F2,F6!} & 4 & 14 & 15 \\
\passthrough{\lstinline!SUB.D F8,F4,F2!} & 5 & 17 & 18 \\
\passthrough{\lstinline!ADD.D F4,F2,F6!} & 8 & 10 & 11 \\
\end{longtable}
}

对比记分牌：Tomasulo 更快，主要因为重命名消除了 WAR/WAW，使后面的
\passthrough{\lstinline!SUB.D!}、\passthrough{\lstinline!ADD.D!}
可以更早流出。

\subsubsection{8.4
作业题考前背法}\label{ux4f5cux4e1aux9898ux8003ux524dux80ccux6cd5}

\begin{enumerate}
\def\labelenumi{\arabic{enumi}.}
\tightlist
\item
  第一次作业：看到``CPI/MIPS/执行时间''，先列
  \passthrough{\lstinline!CPU time = IC×CPI×Cycle Time!}。
\item
  第一次作业：看到多个程序性能平均，问自己是``时间''\,``相对性能''还是``速率''。
\item
  第二次作业：看到流水线时空图，先找瓶颈段和总完成时间。
\item
  第二次作业：看到预约表，按同一行位置差求禁止表。
\item
  第三次作业：看到记分牌，按``结构冲突、WAW、RAW、WAR''逐条卡。
\item
  第三次作业：看到 Tomasulo，记住 \passthrough{\lstinline!Q!}
  是等谁，\passthrough{\lstinline!V!} 是已有值，寄存器状态表存 tag。
\end{enumerate}

\subsection{9.
高频综合题模板}\label{ux9ad8ux9891ux7efcux5408ux9898ux6a21ux677f}

\subsubsection{9.1 性能题模板}\label{ux6027ux80fdux9898ux6a21ux677f}

遇到性能题，按顺序写：

\begin{lstlisting}
CPU time = IC × CPI × Cycle Time
若只比较同一程序且 IC 相同，比较 CPI × Cycle Time
若有局部优化，用 Amdahl
若有停顿，CPI = 理想 CPI + 各类停顿 CPI
\end{lstlisting}

\subsubsection{9.2
流水线题模板}\label{ux6d41ux6c34ux7ebfux9898ux6a21ux677f}

\begin{enumerate}
\def\labelenumi{\arabic{enumi}.}
\tightlist
\item
  写出阶段：IF、ID、EX、MEM、WB。
\item
  标出每条指令读写寄存器。
\item
  判断 RAW/WAR/WAW。
\item
  RAW 优先考虑旁路，不能旁路再停顿。
\item
  分支题用 \passthrough{\lstinline!分支频率 × 失败率 × 惩罚!}。
\end{enumerate}

\subsubsection{9.3 Cache 题模板}\label{cache-ux9898ux6a21ux677f}

\begin{enumerate}
\def\labelenumi{\arabic{enumi}.}
\tightlist
\item
  算块大小得到 offset。
\item
  算行数或组数得到 index。
\item
  剩余位数是 tag。
\item
  访问序列题逐项模拟：先算块号，再算行号/组号，看 tag 是否匹配。
\item
  AMAT 题注意 miss rate 是百分数还是小数。
\end{enumerate}

\subsubsection{9.4
动态调度题模板}\label{ux52a8ux6001ux8c03ux5ea6ux9898ux6a21ux677f}

\begin{enumerate}
\def\labelenumi{\arabic{enumi}.}
\tightlist
\item
  RAW：真相关，必须等待数据。
\item
  WAR/WAW：名相关，可用寄存器重命名消除。
\item
  Tomasulo：看保留站、tag、CDB。
\item
  ROB：看是否按序提交，是否支持精确异常。
\end{enumerate}

\subsection{10.
考前最后一页}\label{ux8003ux524dux6700ux540eux4e00ux9875}

必背公式：

\begin{lstlisting}
CPU time = IC × CPI × Cycle Time
CPI = Σ(比例_i × CPI_i)
Speedup = T_old / T_new
Amdahl = 1 / ((1 - f) + f/S)
流水时间 = (k + n - 1) × Δt
流水加速比 = nk / (k + n - 1)
AMAT = Hit Time + Miss Rate × Miss Penalty
多级 AMAT = HT1 + MR1 × (HT2 + MR2 × MP2)
Availability = MTTF / (MTTF + MTTR)
串联系统失效率 = Σ(1/MTTF_i)
\end{lstlisting}

必背判断：

\begin{lstlisting}
RAW = 写后读，真相关，不能靠重命名消除
WAR = 读后写，名相关，可重命名
WAW = 写后写，名相关，可重命名
直接映像 = 位置唯一，简单快但冲突多
全相联 = 任意位置，冲突少但硬件复杂
组相联 = 折中
写直达 = 同时写 Cache 和主存，简单但流量大
写回 = 只写 Cache，替换时写回，复杂但流量小
写分配 = 写不命中时先把块调入 Cache
不写分配 = 写不命中时直接写下一级
\end{lstlisting}

最后建议：先把第 1、3、5、7、8 章的计算题练熟，再回头背第 2 章 ISA
概念和第 4
章向量机概念。期末卷通常不是考``记住多少名词''，而是考你能不能把公式、冲突类型、Cache
地址划分这三类东西算稳。

\end{document}
